\subsection{Requirements}
\subsubsection{User Requirements}
User Requirements articulate the functional needs of the system from the perspective of the end-users.

\paragraph{Student Requirements}
\begin{itemize}
    \item \textbf{UR-STU-01:} As a Student, I want to consume multimedia course materials and complete spaced-repetition quiz reviews scheduled by the system so that I can build durable retention of lesson content.
    \item \textbf{UR-STU-02:} As a Student, I want to participate in context-aware mock interviews through a unified hybrid experience supporting typed and spoken turns so that I can practise articulating my understanding under realistic conditions.
    \item \textbf{UR-STU-03:} As a Student, I want to receive the binary pass/fail outcome and qualitative remediation guidance for each interview session --- without numeric rubric scores, per-outcome verdicts, or model reasoning --- so that I can understand how to improve while being evaluated on overall demonstrated competence rather than on isolated answers.
    \item \textbf{UR-STU-04:} As a Student, I want to track my own progress through learner-facing metrics --- per-lesson Knowledge Retention Estimate, Progression Readiness, and Review Compliance Rate --- so that I can self-regulate my study pace without depending on instructor intervention.
\end{itemize}

\paragraph{Teacher Requirements}
\begin{itemize}
    \item \textbf{UR-TCH-01:} As a Teacher, I want to upload and organise learning materials so that I can build a structured knowledge base for each lesson.
    \item \textbf{UR-TCH-02:} As a Teacher, I want the system to draft quiz questions and interview prompts from my uploaded materials, and I want to review, edit, approve, or reject each draft before publication, so that AI-generated assessments retain pedagogical accuracy through Human-in-the-Loop oversight.
    \item \textbf{UR-TCH-03:} As a Teacher, I want to confirm the expected response time $T_{\exp}$ for every published quiz card and define the outcome criteria for every published interview so that the system can compute response-time ratios and binary pass/fail verdicts against my domain judgment rather than against generic heuristics.
    \item \textbf{UR-TCH-04:} As a Teacher, I want to configure per-lesson spaced-repetition parameters --- the minimum Easiness Factor for unlock ($EF_{\min}^{\text{unlock}}$), the coverage threshold ($\tau$), and lesson prerequisites --- so that progression stringency can be tuned to the criticality of each lesson.
    \item \textbf{UR-TCH-05:} As a Teacher, I want a cohort monitoring dashboard exposing the class-wide distribution of Knowledge Retention Estimates, a per-student breakdown of an individual's review history, an at-risk list flagging low compliance, frozen retention and a widening gap between recognition and explanation, and the cards the cohort finds hardest, so that I can identify at-risk learners and detect systemic weak points in my source material.
\end{itemize}

\paragraph{Educational Manager Requirements}
\begin{itemize}
    \item \textbf{UR-MGR-01:} As an Educational Manager, I want to enrol Students onto a course or a learning programme in bulk from a CSV file --- with accounts created automatically for addresses the platform has not seen, and per-row failures reported without losing the rest of the batch --- and to invite individual Teachers and Students into my organisation with the role they are to hold, so that I can provision a cohort without creating every account by hand.
    \item \textbf{UR-MGR-02:} As an Educational Manager, I want to define career pathways as versioned sequences of ordered stages, each grouping the courses a Student must satisfy before the next opens, and to assign Teachers to those courses, so that my institution's curriculum structure is reflected in the platform's progression model and a revision published mid-enrolment does not change what an enrolled Student must finish.
    \item \textbf{UR-MGR-03:} As an Educational Manager, I want to view organisation-scoped Career Readiness snapshots --- aggregated across all enrolled Students in a pathway, with drill-down to individual progress --- so that I can report on cohort outcomes without exposing per-interview rubrics.
    \item \textbf{UR-MGR-04:} As an Educational Manager, I want all my actions to be strictly scoped to my single tenant organisation so that I cannot accidentally or intentionally access data belonging to other institutions on the platform.
\end{itemize}

\paragraph{Faculty Dean Requirements}
\begin{itemize}
    \item \textbf{UR-DN-01:} As a Faculty Dean, I want to review, acknowledge, approve, or reject eligible student pathway-change and pathway-drop requests within programmes owned by my faculty, with a recorded reason and no self-approval, so that pathway governance is accountable and scoped to the appropriate academic authority.
\end{itemize}

\paragraph{IT Administrator Requirements}
\begin{itemize}
    \item \textbf{UR-ITA-01:} As an IT Administrator, I want to provision new tenant organisations and invite their initial Educational Manager so that institutions can be onboarded onto the platform without my involvement in their internal user management.
    \item \textbf{UR-ITA-02:} As an IT Administrator, I want to tune the platform's operational parameters at runtime --- enrolment ceilings, retention windows, AI timeout and rate-limit behaviour, ingestion and retrieval thresholds, and the toggles that enable or disable optional processing stages --- with an organisation-level override where a tenant needs to differ, so that the deployment can be adjusted to institutional demand without a code change.
    \item \textbf{UR-ITA-03:} As an IT Administrator, I want to monitor cross-tenant system health, infrastructure metrics, and retained append-only audit records so that I can ensure platform stability and respond to security or compliance incidents.
    \item \textbf{UR-ITA-04:} As an IT Administrator, I want every action I take inside a tenant --- restoring deleted content, disabling an account, amending an organisation's configuration --- to be recorded in retained append-only audit records, so that platform authority broad enough to recover a tenant from an operational failure remains accountable, and pedagogical decisions stay in practice with the Educational Manager and Teachers even where the permission model does not forbid me from making them.
    \item \textbf{UR-ITA-05:} As an IT Administrator, I want the platform to enforce centralised authentication and identity controls --- including Google SSO, two-factor authentication, session management, and role-based access control --- so that user access is uniformly secured across all tenant organisations.

\end{itemize}
\subsubsection{Functional Requirements}

The functional requirements of aBridgeAI specify the concrete system behaviours, processes, and user interactions that the platform must support. They are organised into eight functional modules: (1)~Authentication and Access Control, (2)~Organisation and User Management, (3)~Course and Material Management, (4)~Learning Loop and Adaptive Quiz, (5)~Application Loop and Mock Interview, (6)~Analytics, Reporting, and Notifications, (7)~Learning Programs and Path Governance, and (8)~Platform Configuration and Policy Documents. Each requirement is traced to one or more User Requirements to ensure comprehensive coverage and end-to-end traceability. The complete specification --- one hundred requirements, each with its actors, traced user requirement, feasibility and testability assessment, and priority --- is given in Appendix~\ref{appendix:fr}, and is not reproduced here.

The modules divide along the lines the system itself is built on. Modules~1 and~2 establish who a user is and what they may do: federated sign-in, second-factor verification, and the dynamic role and permission model, then the tenant structure of organisations, faculties and memberships within which every later permission is resolved. Modules~3 to~5 carry the teaching loop in the order a learner meets it --- a course and its ingested materials, the quiz cycle and the spaced-repetition scheduling that governs progression, and the mock interview that tests explanation rather than recognition. Module~6 turns what those produce into reporting, notification and audit.

The final two modules are the ones most often mistaken for administration rather than pedagogy. Module~7 governs the enrolment structure above individual courses --- versioned learning programmes, the career pathways they bundle, and the dean-adjudicated workflow by which a student moves between pathways without losing earned progress. Module~8 covers the runtime configuration catalogue and the versioned policy documents, the two surfaces through which the platform itself is tuned and its terms published.

Two properties hold across the assessment requirements and are worth stating once rather than repeating in each requirement. AI-drafted quiz and interview questions enter a review state and require a human decision before publication; automatically derived ingestion units instead support retrieval and generation. Rules that judge a learner --- such as the integrity policy on an attempt, the pathway version on an enrolment, and the programme version on a student --- are frozen when they start applying, so that editing the curriculum does not retroactively change what somebody already under way is measured against.



% =================================================
% ========================= NFRs ===================
% =================================================
\subsubsection{Non-Functional Requirements}

Non-functional requirements define the quality attributes and operational constraints of aBridgeAI, specifying how the system behaves rather than what it does. They govern performance, scalability, security, and usability properties that collectively determine the reliability and institutional suitability of the platform. The complete Non-Functional Requirements specification is provided in Appendix~\ref{appendix:nfr}. Representative points are summarised below:

\begin{itemize}
    \item Performance and latency targets ensure AI interview responses and core learning workflows remain responsive for users.
    \item Scalability and reliability requirements target modest peak concurrency while prioritising graceful degradation and uptime.
    \item Security and compliance mandates cover encryption at rest, RBAC, audit logging, session management, and the input and output guards that protect the AI interviewer against prompt injection and content leakage.
    \item Usability and accessibility requirements emphasise responsive UI and WCAG 2.1 AA alignment.
\end{itemize}

\noindent Two of these carry more weight than their priority suggests. The concurrency target is modest by design --- the deployment is institutional rather than public, and the measured headroom against it is reported in Section~\ref{sec:evaluation} --- but the accessibility requirement is the one most likely to be treated as optional and least likely to be recoverable late, since retrofitting keyboard navigation and contrast onto a finished component library costs far more than building to them. The nineteen requirements, with their actors and priorities, are enumerated in Appendix~\ref{appendix:nfr}.




\subsubsection{Data Requirements}

This section summarises the data demands of aBridgeAI across its core
functional domains. Rather than specifying a physical schema, the data
requirements identify what information must be captured, maintained,
and related to support each system capability. The complete
specification --- fifteen domains, covering the full lifecycle of
content creation, assessment, learner modeling, and progress reporting
--- is provided in Appendix~\ref{appendix:data}. The governing
demands are as follows.

\begin{itemize}
    \item \textbf{Audit and traceability.} Mutable, auditable domain
    records capture the timestamps and responsible actors needed for
    compliance reporting and dispute resolution; soft-deletable records
    additionally record deletion metadata. Append-only event and history
    records may retain only their creation timestamp and actor context
    appropriate to that event.

    \item \textbf{Identity and access.} One account belongs to at most
    one organisation at a time. Permissions are not hardcoded to roles
    but defined at runtime, and role assignments are scoped to one of
    four levels --- global, organisation, organisational unit, or
    course --- so that the same person may hold different authority in
    different contexts.

    \item \textbf{Curriculum structure.} Organisations contain
    faculties and organisation-owned career pathways. Faculty-scoped
    learning programmes bundle pinned pathway versions; courses divide
    into modules, and each module holds lessons, quizzes, and interviews
    as sibling curriculum items.

    \item \textbf{Content provenance.} Uploaded file content is retained
    through material versions while material metadata remains editable.
    Quiz and interview questions record their generation provenance and
    review state; automatically derived ingestion units support retrieval
    but are not individually teacher-review artefacts.

    \item \textbf{Assessment records.} Spaced-repetition state is held
    per student-card pair in isolation from every peer, while quizzes
    additionally maintain a single defensible grade of record per
    student. Unlock and completion are stored separately, because
    entering a piece of curriculum and discharging it are distinct
    facts.

    \item \textbf{Integrity and safety.} Session activity signals and
    AI guard decisions are recorded as evidence for human review; the
    guard trail is redacted by design, holding a fingerprint of an
    offending input rather than the input itself. The integrity policy
    is frozen onto an attempt when it starts, so a later edit cannot
    re-judge work already under way. Where a quiz requires a camera,
    the current implementation checks only that an active video track
    is present and keeps the stream local to the learner's browser.

    \item \textbf{Operational telemetry.} External AI calls record
    available token usage, latency, and outcome. When token usage and a
    model-price entry are configured, an estimated cost is persisted at
    call time so that later price changes cannot rewrite that estimate.
\end{itemize}
